\begin{table}[h]
\footnotesize
\centering
\begin{tabular}{|p{4cm}|p{12.5cm}|}
\hline
\textbf{Título:} & DISEÑO EN HARDWARE DEL ALGORITMO HAMMING QUASI-CYLIC (HQC) PARA ENCRIPTACIÓN POSTCUÁNTICA \\ \hline

\textbf{Palabras clave:} & Criptografía postcuántica, computación cuántica, diseño en FPGA, encapsulación de clave, ciberseguridad. \\ \hline

\textbf{Resumen:} & 
\begin{minipage}[t]{\linewidth}
\justifying
Este proyecto presenta el diseño en hardware del algoritmo Hamming Quasi-Cyclic, un mecanismo de encapsulación de clave resistente a ataques cuánticos seleccionado por el NIST como quinto algoritmo dentro del proceso de estandarización de criptografía post-cuántica. Debido al avance acelerado de la computación cuántica y su potencial para vulnerar sistemas criptográficos clásicos, se propone evaluar la viabilidad de HQC en plataformas FPGA, optimizando su rendimiento en velocidad, uso de recursos y consumo energético. Para ello, se desarrolla una arquitectura digital en VHDL, validada mediante simulaciones y comparaciones con diseños existentes, aportando al avance de soluciones seguras y eficientes frente a futuras amenazas en ciberseguridad, sentando bases para futuras aplicaciones en sistemas embebidos y redes seguras.
\end{minipage} \\ \hline

\textbf{Área temática:} & Implementación de Algoritmos Criptográficos en Hardware \\ \hline

\textbf{Modalidad de trabajo de grado:} &
$\boxtimes$ Profesional \quad $\Box$ Revisión crítica \quad $\Box$ Investigación e innovación \\
& $\Box$ En la industria/empresa \quad $\Box$ Creación de empresa \quad $\Box$ Práctico social \\ \hline

\textbf{Duración del trabajo de grado:} & 10 Meses \\ \hline
\textbf{Entidades participantes:} & Universidad del Valle \\ \hline

\textbf{Costo presupuestado del trabajo:} &
\makecell[tl]{%
Aportes en especie: COP\$ 18'672.000\\
Aportes en efectivo: COP\$ 0\\
Costo global: COP\$ 18'672.000
} \\ \hline

\begin{comment}
 \textbf{Áreas académicas EIEE:} &
\makecell[tl]{%
$\Box$ Arquitecturas Digitales\\
$\Box$ Bionanoelectrónica\\
$\Box$ Conversión de energía\\
$\Box$ Informática industrial\\
$\Box$ Sistema de control y accionamientos\\
$\Box$ Sistemas de potencia\\
$\Box$ Telecomunicaciones
} \\ \hline   
\end{comment}



\textbf{Grupos de investigación:} &
\makecell[tl]{%
$\Box$ Alta tensión (GRALTA)\\
$\Box$ Arquitecturas digitales y microelectrónica (GADYM)\\
$\boxtimes$ Bionanoelectrónica (BIONANO)\\
$\Box$ Control industrial (GICI)\\
$\Box$ Conversión de energía (CONVERGIA)\\
$\Box$ Percepción y sistemas inteligentes (PSI)\\
$\Box$ Sistemas de telecomunicaciones (SISTEL-UV)\\
$\Box$ Sistema de hidroeléctricos de generación (SIIG)\\
$\Box$ Otro:
} \\ \hline



\multicolumn{2}{|c|}{\textbf{Estudiante}} \\ \hline
\textbf{Nombre:} & Camilo Home Buriticá \hfill \textbf{Código:} 202043403 \\ \hline
\textbf{Correo electrónico:} & camilo.home@correounivalle.edu.co \hfill \textbf{Tel.:} (57) 3226928217 \\ \hline

\multicolumn{2}{|c|}{\textbf{Director}} \\ \hline
\textbf{Nombre:} & Jaime Velasco Medina \hfill \textbf{Tel.:} (57) 6023212100 ext. 3283/7390\\ \hline
\textbf{Correo electrónico:} & jaime.velasco@correounivalle.edu.co  \\ \hline
\textbf{Institución:} & Universidad del Valle - Bionanoelectrónica (BIONANO) \\ \hline


   



\end{tabular}
\end{table}
